\chapter{Poisson Distribution}
\label{chap:poisson-distribution}
\lecture{3}{15 May. 11:00}{Third Lecture}

\begin{eg}[Simpler Example: Customers at a Cafe]
  \label{ex:cafe-customers}
  Imagine a small cafe where, on average, 2 customers arrive every minute during the lunch rush. We want to find the probability that exactly 5 customers will arrive in a specific one-minute interval during this period. This scenario can be modeled using the Poisson distribution.

  Here, the factors of the Poisson distribution are:
  \begin{itemize}
    \item \(\lambda\) (lambda): The average rate of arrivals. In this case, \(\lambda = 2\) customers per minute. This average is assumed to be constant during the lunch rush period.
    \item \(x\): The specific number of arrivals we are interested in. Here, \(x = 5\) customers.
    \item \(e\): Euler's number, approximately 2.71828.
  \end{itemize}

  \textbf{Explanation of "specific one-minute interval":}
  The term "specific one-minute interval" refers to any chosen single minute
  within the timeframe where the average arrival rate (\(\lambda=2\)
  customers/minute) is valid – in this case, during the lunch rush. It
  doesn't mean a pre-selected or unique minute, but rather \emph{any} individual minute you decide to observe. For example:
  \begin{itemize}
    \item The minute from 12:05 PM to 12:06 PM.
    \item The minute from 1:10 PM to 1:11 PM.
    \item Any other one-minute slot during the lunch rush.
  \end{itemize}
  The Poisson distribution calculates the probability of observing \(x\) events (5 customers) in such a defined interval, assuming the rate \(\lambda\) is constant across all such intervals and events occur independently. So, for any given minute you pick during the lunch rush, the probability of 5 customers arriving in \emph{that particular minute} is what we are calculating.

  The probability mass function (PMF) for the Poisson distribution is:
  \[ P(X=x) = \frac{e^{-\lambda} \lambda^x}{x!} \]
  We want to calculate \(P(X=5)\) when \(\lambda=2\).

  Plugging the values into the formula:
  \[ P(X=5) = \frac{e^{-2} \cdot 2^5}{5!} \]

  Let's break down the calculation:
  \begin{itemize}
    \item \(e^{-2} \approx 0.135335\)
    \item \(2^5 = 32\)
    \item \(5! = 5 \times 4 \times 3 \times 2 \times 1 = 120\)
  \end{itemize}

  So,
  \[ P(X=5) = \frac{0.135335 \times 32}{120} \]
  \[ P(X=5) = \frac{4.33072}{120} \]
  \[ P(X=5) \approx 0.036089 \]

  Therefore, the probability of exactly 5 customers arriving in any specific one-minute interval during the lunch rush is approximately 0.0361, or about 3.61%. This demonstrates how the Poisson distribution uses the average rate (\(\lambda\)) to predict the likelihood of a specific number of events (\(x\)) occurring in a fixed interval of time[19].
\end{eg}

\section{Introduction}
It is used to model the probability of a given number of events
occurring within a fixed interval of time or space, assuming these events happen
with a known constant mean rate and independently of the time since the last
event.

\begin{definition}[Poisson Distribution]
  \label{def:poisson}
  A discrete random variable \(X\) is said to follow a Poisson distribution if its probability mass function (PMF) is given by:
  \[ P(X=x) = \frac{e^{-\lambda} \lambda^x}{x!} \]
  for \(x = 0, 1, 2, \ldots\)\marginpar{\(\lambda > 0\)}.
  Here:
  \begin{itemize}
    \item \(x\) is the number of occurrences of an event (a non-negative integer).
    \item \(\lambda\) (lambda) is the average number of events occurring in the given interval. It is the sole parameter of the distribution.
    \item \(e\) is Euler's number (approximately 2.71828).
    \item \(x!\) is the factorial of \(x\).
  \end{itemize}
  The notation for a random variable \(X\) following a Poisson distribution with mean \(\lambda\) is \(X \sim \mathcal{P}(\lambda)\) or \(X \sim \text{Poisson}(\lambda)\).
  Conditions for using Poisson distribution include:
  \begin{itemize}
    \item Events occur at random and independently. The occurrence of one event does not affect the probability of another event.
    \item The mean number of events (\(\lambda\)) occurring within a given interval is known and constant.
  \end{itemize}
\end{definition}

\section{Meaning of \(P(X=x)\)}
The expression \(P(X=x)\) represents the probability that exactly \(x\) events will occur in a specified interval of time or space. For instance, if \(X\) represents the number of emails received per hour and \(X \sim \mathcal{P}(10)\), meaning on average 10 emails are received per hour, then \(P(X=5)\) would be the probability of receiving exactly 5 emails in a particular hour. This is calculated using the Poisson formula:
\[ P(X=5) = \frac{e^{-10} 10^5}{5!} \]
The value of \(x\) can be any non-negative integer (0, 1, 2, ...), reflecting that there's theoretically no upper limit to the number of events that can occur.

\section{Derivation of Expectation and Variance}
For a Poisson-distributed random variable \(X \sim \mathcal{P}(\lambda)\), the expectation (mean) and variance are both equal to \(\lambda\). \marginpar{Mean = Variance = \(\lambda\)}

\subsection{Expectation \(E[X]\)}
The expectation of a discrete random variable \(X\) is given by \(E[X] = \sum_{x=0}^{\infty} x P(X=x)\).
\begin{align*}
  E[X] &= \sum_{x=0}^{\infty} x \frac{e^{-\lambda} \lambda^x}{x!} \\
       &= e^{-\lambda} \sum_{x=1}^{\infty} x \frac{\lambda^x}{x!} \quad (\text{The term for } x=0 \text{ is } 0) \\
       &= e^{-\lambda} \sum_{x=1}^{\infty} \frac{\lambda^x}{(x-1)!} \\
  \text{Let } k = x-1, \text{ so } x = k+1. &\text{ When } x=1, k=0. \\
  E[X] &= e^{-\lambda} \sum_{k=0}^{\infty} \frac{\lambda^{k+1}}{k!} \\
       &= e^{-\lambda} \lambda \sum_{k=0}^{\infty} \frac{\lambda^k}{k!} \\
  \text{Since } \sum_{k=0}^{\infty} \frac{\lambda^k}{k!} = e^{\lambda} &\text{ (Taylor series expansion of } e^{\lambda}), \\
  E[X] &= e^{-\lambda} \lambda e^{\lambda} \\
       &= \lambda \label{eq:poisson_expectation}
\end{align*}
Thus, the expectation of a Poisson distribution is \(\lambda\).

\subsection{Variance \(\text{Var}(X)\)}
The variance of a random variable \(X\) is given by \(\text{Var}(X) = E[X^2] - (E[X])^2\).
First, we calculate \(E[X(X-1)]\):
\begin{align*}
  E[X(X-1)] &= \sum_{x=0}^{\infty} x(x-1) \frac{e^{-\lambda} \lambda^x}{x!}  \\
            &= e^{-\lambda} \sum_{x=2}^{\infty} x(x-1) \frac{\lambda^x}{x!} \quad (\text{Terms for } x=0, 1 \text{ are } 0) \\
            &= e^{-\lambda} \sum_{x=2}^{\infty} \frac{\lambda^x}{(x-2)!} \\
  \text{Let } j = x-2, \text{ so } x = j+2. &\text{ When } x=2, j=0. \\
  E[X(X-1)] &= e^{-\lambda} \sum_{j=0}^{\infty} \frac{\lambda^{j+2}}{j!} \\
            &= e^{-\lambda} \lambda^2 \sum_{j=0}^{\infty} \frac{\lambda^j}{j!}  \\
            \text{Since } \sum_{j=0}^{\infty} \frac{\lambda^j}{j!} = e^{\lambda}, \\
  E[X(X-1)] &= e^{-\lambda} \lambda^2 e^{\lambda} \\
            &= \lambda^2 
\end{align*}
We know that \(E[X(X-1)] = E[X^2 - X] = E[X^2] - E[X]\).
So, \(E[X^2] = E[X(X-1)] + E[X]\).
Using \(E[X] = \lambda\) and \(E[X(X-1)] = \lambda^2\):
\[ E[X^2] = \lambda^2 + \lambda \label{eq:poisson_ex2} \]
Now, we can find the variance:
\begin{align*}
  \text{Var}(X) &= E[X^2] - (E[X])^2 \\
                &= (\lambda^2 + \lambda) - (\lambda)^2 \\
                &= \lambda^2 + \lambda - \lambda^2 \\
                &= \lambda \label{eq:poisson_variance}
\end{align*}
Thus, the variance of a Poisson distribution is also \(\lambda\). \marginpar{Standard deviation \(\sigma = \sqrt{\lambda}\)}

\section{Applications of Poisson Distribution}
The Poisson distribution is a versatile tool used in various fields to model the number of occurrences of an event over a specific interval. Some key applications include:
\begin{itemize}
  \item \textbf{Telecommunications:} Predicting the number of calls arriving at a call center per hour, or the number of network failures per week. This helps in staffing and resource allocation.
  \item \textbf{Quality Control:} Modeling the number of defects or errors in a manufactured product (e.g., defects per square meter of material or per batch of items).
  \item \textbf{Finance and Insurance:} Estimating the number of insurance claims per month or bankruptcies filed per month.
  \item \textbf{Biology and Medicine:} Counting the number of bacteria in a culture, radioactive decay particles per minute, or instances of a rare disease in a population. \href{https://www.studypug.com/statistics-help/poisson-distribution}{Ecology example from StudyPug} could involve modeling the number of trees affected by a disease per acre.
  \item \textbf{Astronomy:} Modeling the number of meteorites of a certain size striking the Earth in a year.
  \item \textbf{Traffic Engineering:} Predicting the number of cars arriving at an intersection or a toll booth in a specific time interval.
  \item \textbf{Retail Management:} Forecasting customer arrivals at a store or checkout queues, helping in optimizing staff schedules and queue management.
\end{itemize}
The Poisson distribution is particularly useful for modeling rare events or counting processes where events occur independently and at a constant average rate.
